\documentclass[a4paper,9pt]{extarticle}

%=========================================================
% PACKAGES
%=========================================================

\usepackage[
a4paper,
top=1.8cm,
bottom=1.5cm,
footskip=7mm,
left=1cm,
right=1cm,
headheight=14pt,
headsep=5mm
]{geometry}

\usepackage[T1]{fontenc}
\usepackage[utf8]{inputenc}
\usepackage[french]{babel}

\usepackage[sfdefault]{FiraSans}
\renewcommand{\familydefault}{\sfdefault}

\usepackage{multicol}
\setlength{\columnsep}{0.7cm}

\usepackage{amsmath,amssymb}
\usepackage{titlesec}
\usepackage{enumitem}

\usepackage{tcolorbox}
\tcbuselibrary{skins,breakable}

\usepackage{xcolor}
\usepackage{graphicx}
\usepackage{float}

\usepackage{fancyhdr}
\usepackage[hidelinks]{hyperref}


%=========================================================
% PAGE SETTINGS
%=========================================================

\setlength{\parindent}{0pt}
\setlength{\parskip}{1.5pt}

\setlength{\abovedisplayskip}{4pt}
\setlength{\belowdisplayskip}{4pt}

\setlist[itemize]{
leftmargin=*,
itemsep=1pt,
topsep=1pt
}
\setlist[enumerate]{
leftmargin=*,
itemsep=1pt,
topsep=1pt
}


%=========================================================
% COLORS
%=========================================================

\definecolor{Blue}{HTML}{E3F0FF}
\definecolor{Green}{HTML}{E3F7EA}
\definecolor{Pink}{HTML}{FFE5ED}
\definecolor{Orange}{HTML}{FFF0DC}
\definecolor{Purple}{HTML}{EEE7FF}
\definecolor{Yellow}{HTML}{FFFADB}
\definecolor{Goat}{HTML}{F7EEDF}

\definecolor{BlueTitle}{HTML}{285B8F}
\definecolor{GreenTitle}{HTML}{287052}
\definecolor{PinkTitle}{HTML}{A84063}
\definecolor{OrangeTitle}{HTML}{A96500}
\definecolor{PurpleTitle}{HTML}{6245A5}
\definecolor{GoatTitle}{HTML}{8A5A1F}
\definecolor{RecapTitle}{HTML}{5B3A8E}


%=========================================================
% HEADER
%=========================================================

\pagestyle{fancy}
\fancyhf{}

\fancyhead[L]{\small\textbf{Comportement mécanique des matériaux}}
\fancyhead[C]{\small Al--Cu : solidification, maturation, durcissement structural}
\fancyhead[R]{\small Week 08 --- 26/10-01/11/2026}

\fancyfoot[C]{\small\href{https://github.com/AlphaXT38G/notes_ba5}{github.com/AlphaXT38G/notes\_ba5}}
\fancyfoot[R]{\small \thepage}

\renewcommand{\headrulewidth}{0.4pt}


%=========================================================
% TITLES
%=========================================================

\titleformat{\section}
{\large\bfseries\color{BlueTitle}}
{}{0pt}{}

\titlespacing{\section}
{0pt}{6pt}{3pt}


%=========================================================
% BOXES
%=========================================================
% Boxes are kept whole (unbreakable): inside multicols, a box split by tcolorbox
% leaves an empty column behind it.

\newtcolorbox{definition}{
enhanced, unbreakable, colback=Blue, colframe=BlueTitle,
boxrule=0.7pt, arc=0mm, title=\textbf{Définition}, fonttitle=\bfseries
}

\newtcolorbox{important}{
enhanced, unbreakable, colback=Pink, colframe=PinkTitle,
boxrule=0.7pt, arc=0mm, title=\textbf{Idée clé}, fonttitle=\bfseries
}

\newtcolorbox{formula}{
enhanced, unbreakable, colback=Green, colframe=GreenTitle,
boxrule=0.7pt, arc=0mm, title=\textbf{Formules clés}, fonttitle=\bfseries
}

\newtcolorbox{warning}{
enhanced, unbreakable, colback=Orange, colframe=OrangeTitle,
boxrule=0.7pt, arc=0mm, title=\textbf{Attention}, fonttitle=\bfseries
}

\newtcolorbox{example}{
enhanced, unbreakable, colback=Purple, colframe=PurpleTitle,
boxrule=0.7pt, arc=0mm, title=\textbf{Exemple}, fonttitle=\bfseries
}

% GOAT - brown: the idea told as simply as possible
\newtcolorbox{goat}{
enhanced, unbreakable, colback=Goat, colframe=GoatTitle,
boxrule=0.7pt, arc=0mm, title=\textbf{Expliqué à une chèvre}, fonttitle=\bfseries
}

% RECAP - violet: what to know for the exam, at the end of the week
\newtcolorbox{recap}{
enhanced, breakable, colback=Purple, colframe=RecapTitle,
boxrule=1pt, arc=0mm, title=\textbf{Récap pour l'examen}, fonttitle=\bfseries
}

% Screenshot of a slide of the course: \slide{cNN-PP}
\newcommand{\slide}[2][0.95]{%
\begin{center}
\includegraphics[width=#1\linewidth]{../images/course/#2}
\end{center}}


%=========================================================
% DOCUMENT
%=========================================================

\begin{document}

\thispagestyle{fancy}

\begin{multicols}{2}

%=========================================================
\section{Solidification de Al--Cu}

\begin{definition}
Côté Al : eutectique $\alpha$ (Cu dans Al, max 5,7\,\% à 548\,°C) + $\theta$ (Al$_2$Cu,
fragile) à 33,2\,\% Cu. Points triples : eutectique, eutectoïde, péritectique, péritectoïde.
\end{definition}

\begin{formula}
Levier : $f_\alpha = \dfrac{C_\beta - C}{C_\beta - C_\alpha}$, $f_\beta = \dfrac{C - C_\alpha}{C_\beta - C_\alpha}$
(bras \textbf{opposé}).
\end{formula}

\begin{important}
Al--4,5Cu : dendrites $\alpha$. Al--33Cu : eutectique (croissance \textbf{couplée}). Al--10Cu :
dendrites puis eutectique ($\frac{10-5{,}7}{33{,}2-5{,}7} = 16\,\%$ d'eutectique, \textit{calcul pas dans
les diapos}).
\end{important}

\slide[0.8]{c03b-16}

\begin{warning}
Le levier suppose l'équilibre partout ; dans le solide la diffusion est trop lente : le solide
garde sa composition $\Rightarrow$ \textbf{microségrégation} (gradients, $\theta$ hors équilibre
dans Al--4,5Cu brut de coulée). Sauf soluté interstitiel (C dans Fe).
\end{warning}

\begin{goat}
Les premiers cristaux (des sapins) sont pauvres en cuivre, le dernier liquide très riche gèle en
lamelles, et comme rien ne bouge dans le solide, ça reste figé.
\end{goat}

%=========================================================
\section{Maturation et homogénéisation}

\begin{formula}
\[ \Delta P = \frac{2\gamma}{R}, \quad \Delta G = \frac{2\gamma V_m}{R}, \quad \lambda^3 - \lambda_0^3 = Kt \ (l = kt^{1/3}) \]
\end{formula}

\begin{important}
Petites particules moins stables (deux ballons : le petit se vide dans le gros) : la structure
grossit. Refroidir vite $\Rightarrow$ plus fin ; très vite : phases hors équilibre, voire amorphe.
\end{important}

\slide[0.85]{c03b-23}

\begin{warning}
\textbf{Homogénéisation} (dissoudre $\theta$, égaliser) \textbf{sous l'eutectique} : sinon
l'eutectique fond, mûrit, $\theta$ aux joints : alliage \textbf{brûlé}, fragile.
\end{warning}

\begin{goat}
Les grosses gouttes mangent les petites ; la structure grossit comme la racine cubique du temps.
\end{goat}

%=========================================================
\section{Cinétique et TTT}

\begin{important}
Thermodynamique : \textbf{si} ($\Delta G$). Cinétique : \textbf{à quelle vitesse}. Lente près de
l'équilibre (pas de force motrice) et à froid (pas de diffusion) $\Rightarrow$ \textbf{nez}.
Isotherme : courbe en S.
\end{important}

\slide[0.75]{c03b-35}

\begin{definition}
\textbf{TTT} : trempe infiniment rapide à $T$ puis maintien ; début et fin de transformation.
Valable pour des histoires isothermes.
\end{definition}

\begin{goat}
Il faut à la fois en avoir envie (froid) et pouvoir bouger (chaud) : le plus rapide est au milieu.
\end{goat}

%=========================================================
\section{Durcissement structural}

\begin{definition}
\textbf{Mise en solution} ($\alpha$ seul) $\to$ \textbf{trempe} (sursaturation : à l'ambiante
$\approx 0{,}1\,\%$ Cu toléré) $\to$ \textbf{revenu} (à chaud) ou \textbf{vieillissement}
(ambiante).
\end{definition}

\slide[0.8]{c03b-42}

\begin{important}
Séquence GP $\to \theta'' \to \theta' \to \theta$. Interfaces \textbf{cohérentes} (faible
$\gamma$, champ de contraintes), semi-cohérentes (dislocations), incohérentes ($\gamma$
élevée). Les cohérents germent plus facilement : ils viennent d'abord. Moins stable = plus basse
$T$ max de formation, plus de Cu requis.
\end{important}

\begin{formula}
Contournement (Orowan) : $\sigma = \dfrac{2Gb}{d} \approx \dfrac{Ea}{d}$ ($d$ : distance entre
précipités).
\end{formula}

\slide[0.75]{c03b-48}

\begin{important}
Petits précipités \textbf{cisaillés} (durcit avec la taille) ; gros \textbf{contournés}, et la
maturation les écarte ($d$ $\uparrow$) $\Rightarrow$ \textbf{optimum}, puis sur-revenu. Dépend
de la température.
\end{important}

\begin{warning}
Conditions : solubilité décroissante avec $T$ ; cinétique praticable ; précipités efficaces
(pas Al--Mg) ; pas d'usage là où la maturation agit.
\end{warning}

\begin{goat}
On piège le cuivre, puis il ressort en poussière de grains où les dislocations se cognent.
Trop attendre : les grains grossissent, s'écartent, et le métal ramollit.
\end{goat}

%=========================================================
\section{Récap de la semaine}

\begin{recap}
\begin{itemize}
\item Levier (bras opposé) ; 3 alliages Al--Cu ; dendrites, eutectique.
\item Microségrégation : solide figé, phases hors équilibre.
\item $\Delta P = 2\gamma/R$ ; $\lambda^3 - \lambda_0^3 = Kt$ ; homogénéiser sous l'eutectique.
\item TTT : nez = force motrice vs diffusion.
\item Durcissement structural : solution, trempe, revenu ; GP $\to\theta$ ; cohérence ;
      cisaillement puis $\sigma = 2Gb/d$ ; optimum ; conditions.
\end{itemize}
\end{recap}

\end{multicols}

\end{document}
